\section{Conclusion}
\label{sec:conclusion}

Les travaux examinés fournissent les pièces du problème, dispersées entre plusieurs domaines. La fouille d’arguments montre que les relations sont plus difficiles à retrouver que les composantes, et que la structure à longue distance échappe à une analyse phrase par phrase. Les systèmes qui visualisent des extractions de LLM montrent qu’une extraction n’est utilisable qu’avec des boucles de contrôle, et que les utilisateurs gardent leur confiance même lorsqu’ils voient des erreurs. Les interfaces de lecture montrent comment augmenter un document sans le remplacer. Les travaux sur les citations et sur les LLM employés comme juges rappellent qu’une citation n’est pas une preuve et qu’un modèle ne peut pas être le seul juge de ce qu’il produit.

Visual Reasoning Maps réunit ces leçons dans une même interface : une carte de taille limitée, des relations typées, un renvoi vérifié à la phrase et un statut visible pour chaque élément. Reste la question que seul le protocole d’évaluation peut trancher : la carte aide-t-elle le lecteur à juger un raisonnement, ou lui donne-t-elle seulement l’impression de l’avoir compris ?
